% Frame 20: the distribution mu* of the barrier (Theorem 7.1) on the classes A, D, T_1, T_2.
% Adapted from ../paper/figures/barrier.tex (Figure 5 of the paper).  Mathematical content unchanged:
% same boxes (x-units as in the paper: A=[0,1], D=[1.75,3.75], T_1=[4.5,7.5], T_2=[8.25,11.25]), one dot
% per tuple at the same positions with area proportional to its mu*-mass (5/87 on A, 1/29 = 3/87
% elsewhere, radii in the ratio sqrt(5/3) as in the paper), same class masses, same cells (sat_t shaded;
% conditional probabilities 9/18, 9/18 and 5/20, 6/20, 9/20), same right-hand column.
% Changes for the slide (layout only): larger fonts (the fractions in the cells are scaled by 1.15),
% x-unit 0.82cm, lower boxes and taller cells; the mass label of the narrow box A is moved 2.4mm to the
% left; in the right-hand column the certificate's guarantee (8/9, 9/10) is set in the accent colour and
% the result 27/29 in bold; the description of mu* (Section 7 of the paper: uniform within each class,
% class masses 5:6:9:9) is added at the top right.
% The legend of the cells (caption of Figure 5) is given by the three text zones under the figure on
% the slide.  The data lines of the \foreach lists below were compared with the paper's file by a script.
\begingroup
\def\pVrh{0.32}\def\pVbt{1.16}% row height of the dots; top of the boxes (3*0.32 + 2*0.1)
\def\pVra{1.3mm}%      dot radius on A (mass 5/87)
\def\pVrb{1.00698mm}% dot radius elsewhere (mass 1/29 = 3/87) = 1.3mm*sqrt(3/5); area ratio 5:3
\def\pVch{0.6}\def\pVcp{0.69}% height and pitch of the cells
\begin{tikzpicture}[x=0.82cm,y=1cm,
  cls/.style={draw,rounded corners=2.5pt,line width=0.7pt,fill=black!3},
  cname/.style={font=\Large,inner sep=1pt,text height=1.7ex,text depth=0.5ex},
  cmass/.style={font=\normalsize,inner sep=1pt},
  rlab/.style={font=\normalsize,anchor=east,inner sep=1pt},
  cell/.style={draw=black!45,line width=0.5pt},
  hi/.style={cell,fill=black!12},
  val/.style={font=\normalsize,scale=1.15},
  satv/.style={font=\normalsize},
  bnd/.style={font=\normalsize,anchor=west,inner sep=1pt,scale=1.12}]
% ---------------------------------------------------------------- boxes and dots
\foreach \l/\r in {0/1, 1.75/3.75, 4.5/7.5, 8.25/11.25}
  \draw[cls] (\l,0) rectangle (\r,\pVbt);
\foreach \i in {0,1,2}{
  \fill[black!60] (0.5,\pVbt-0.1-\i*\pVrh-0.5*\pVrh) circle[radius=\pVra];
  \foreach \x in {2.25,3.25,5,6,7,8.75,9.75,10.75}
    \fill[black!60] (\x,\pVbt-0.1-\i*\pVrh-0.5*\pVrh) circle[radius=\pVrb];}
\node[rlab] at (-0.25,0.5*\pVbt) {$\mu^\ast$};
% ---------------------------------------------------------------- headers
\foreach \c/\n/\m in {0.5/{\mathcal A}/{\tfrac{5}{29}=3\cdot\tfrac{5}{87}},
                      2.75/{\mathcal D}/{\tfrac{6}{29}=6\cdot\tfrac{1}{29}},
                      6/{\mathcal T_1}/{\tfrac{9}{29}=9\cdot\tfrac{1}{29}},
                      9.75/{\mathcal T_2}/{\tfrac{9}{29}=9\cdot\tfrac{1}{29}}}{
  \node[cname] at (\c,\pVbt+0.75) {$\n$};
  \ifdim\c pt<1pt% the label of the narrow box A is moved 2.4mm to the left, away from the label of D
    \node[cmass,xshift=-2.4mm] at (\c,\pVbt+0.29) {$\m$};
  \else
    \node[cmass] at (\c,\pVbt+0.29) {$\m$};
  \fi}
% ---------------------------------------------------------------- description of mu* (top right)
\node[anchor=north west,font=\normalsize,inner sep=1pt,align=left] at (11.55,\pVbt+0.96)
  {$\mu^\ast$: \textbf{uniform} within\\[0.3mm]each class, masses\\[1.1mm]\hspace*{4mm}{\large$\mathbf{5:6:9:9}$}};
% ---------------------------------------------------------------- certificate grid
\node[font=\normalsize,anchor=south west,inner sep=1pt] at (11.55,-0.17) {$\mathbb{E}_{\mu^\ast}V_t\,/\,C_t\ \ge$};
\foreach \t/\a/\d/\p/\q/\bd in {
   0/sat/sat/{\tfrac{9}{18}}/{\tfrac{9}{18}}/{\tfrac{11}{29}+\tfrac{18}{29}\cdot{\color{acc}\boldsymbol{\tfrac{8}{9}}}=\boldsymbol{\tfrac{27}{29}}},
   1/{\tfrac{5}{20}}/{\tfrac{6}{20}}/sat/{\tfrac{9}{20}}/{\tfrac{9}{29}+\tfrac{20}{29}\cdot{\color{acc}\boldsymbol{\tfrac{9}{10}}}=\boldsymbol{\tfrac{27}{29}}},
   2/{\tfrac{5}{20}}/{\tfrac{6}{20}}/{\tfrac{9}{20}}/sat/{\tfrac{9}{29}+\tfrac{20}{29}\cdot{\color{acc}\boldsymbol{\tfrac{9}{10}}}=\boldsymbol{\tfrac{27}{29}}}}{
  \pgfmathsetmacro\yt{-0.23-\t*\pVcp}
  \pgfmathsetmacro\yb{\yt-\pVch}
  \node[rlab] at (-0.25,{(\yt+\yb)/2}) {$t=\t$};
  \foreach \l/\r/\v in {0/1/\a, 1.75/3.75/\d, 4.5/7.5/\p, 8.25/11.25/\q}{
    \edef\tmp{\v}\def\satstr{sat}%
    \ifx\tmp\satstr
      \draw[hi] (\l,\yt) rectangle node[satv] {$\sat_{\t}$} (\r,\yb);
    \else
      \draw[cell] (\l,\yt) rectangle node[val] {$\v$} (\r,\yb);
    \fi}
  \node[bnd] at (11.55,{(\yt+\yb)/2}) {$\bd$};}
\end{tikzpicture}
\endgroup
